\documentclass[a5paper,10pt]{article}

% https://github.com/InDevRus/filippov-solutions

% Нумерация страниц
\pagenumbering{gobble}

% Настройка полей
\usepackage{geometry}
\geometry{left=1cm}
\geometry{right=1cm}
\geometry{top=1cm}
\geometry{bottom=1.5cm}

% Вставка таблицы
\usepackage{array}
\usepackage{tabularx}

% Инструменты выравнивания формул
\usepackage{amsmath}
\usepackage{mathtools}

% Диакритика в математике
\usepackage{amsfonts}

% Вычисления размеров на ходу
\usepackage{calc}

% Удобные записи производных
\usepackage[italic=true]{derivative}

% Настройки сносок
\usepackage{enotez}

% Длинные знаки векторов
\usepackage{anyfontsize}
\usepackage[b]{esvect}

% Вставка картинок
\usepackage{graphicx}

% Вставка красной строки
\usepackage{indentfirst}
\setlength{\parindent}{1em}

% Условия в командах
\usepackage{xifthen}

% Луа-скрипты
\usepackage{luacode}

% Настройка команд отдельно в математическом и текстовом режимах
\usepackage{mathcommand}

% Для повторения LaTeX кода
\usepackage{multido}

% Конфигурация отступов формул
\delimitershortfall-1sp
\usepackage{mleftright}
\mleftright

% Растягивание объектов
\usepackage{relsize}

% Обтекание текстом
\usepackage{wrapfig}
\usepackage{subcaption}
\usepackage{floatflt}

% Поддержка ввода кириллицы
\usepackage[english, russian]{babel}

% Настройка корневой позиции для компилятора
\usepackage{import}
\providecommand*{\ProjectRootPath}{..}

\import{\ProjectRootPath/preambles/macros/}{theorems}

\import{\ProjectRootPath/preambles/fonts}{font_setup}

% Знак градуса
\usepackage{siunitx}

\import{\ProjectRootPath/preambles/macros/}{operators}
\import{\ProjectRootPath/preambles/macros/}{redefinitions}

\import{\ProjectRootPath/preambles/fonts}{letters_setup}
\import{\ProjectRootPath/preambles/fonts/adjustments}{custom_superscripts}

\import{\ProjectRootPath/preambles/custom}{formatting}
\import{\ProjectRootPath/preambles/custom}{frames}
\import{\ProjectRootPath/preambles/custom}{list_setup}

% TODO: перевести команды на современный стиль объявления

\begin{document}

Если особая точка уравнения $\br{ax + by} \di x + \br{mx + ny} \di y = 0$ является <<центром>>,
то корни соответствующего характеристического уравнения чисто мнимы.

Не считая тех областей, где $mx + ny = 0$, $\der{y}{x} = \dfrac {-ax - by} {mx + ny}$.

$A = \begin{pmatrix} m & n \\ -a & -b \end{pmatrix}$,
$\br{\lambda - m}\br{\lambda + b} + an = 0$.
$\pow{\lambda}{2} + \br{b - m}\lambda + an - bm = 0$. Если корни комплексные, 
то $\Re\br{\lambda} = \dfrac {b - m} {2} = 0$,
так что $m = b$, и исходное уравнение принимает вид 
$\br{ax + by} \di x + \br{bx + ny} \di y = 0$.
\begin{align*}
    P\br{x; y} = ax + by &&
    Q\br{x; y} = bx + ny &&
    \parder{Q}{x} - \parder{P}{y} = b - b = 0
\end{align*}
Следовательно, данное уравнение -- это уравнение в полных дифференциалах, и его общее решение
выражается квадратичной формой вида
$\dfrac {a}{2} \pow{x}{2} + bxy + \dfrac {n} {2} \pow{y}{2} = C$.

Обратное неверно, так как равенство $m = b$ никак не гарантирует факт, 
что корни характеристического уравнения комплексные. Возможна ситуация, что 
$an - \pow{b}{2} < 0$, и тогда особая точка будет <<седлом>>. Если же $an - \pow{b}{2} = 0$, 
то уравнение примет вид $a \di x + b \di y = 0$, и начало координат вообще не будет
особой точкой.

\end{document}
